\hypertarget{exercises}{}
\hypertarget{exercises-with-guided-solutions}{%
\chapter{Exercises with guided
solutions}\label{exercises-with-guided-solutions}}

\textbf{Exercise 1 --- Unit conversion}

A monthly alpha is 0.00085. The linear annualised value is 0.00085 \ensuremath{\times } 12
\ensuremath{\times } 10,000 = \textbf{102 basis points/year}. This does not compound
monthly returns.

\textbf{Exercise 2 --- Why is a many-to-many merge dangerous?}

It replicates the same economic observation. Point estimates may be
reweighted and standard errors treat duplicate rows as new information.
Validate key uniqueness before and after every join.

\textbf{Exercise 3 --- Joint rejection with mostly uncertain individual ranks}

A joint test aggregates evidence across the alpha vector using
covariance. It can reject all-zero even while many individual posterior
intervals cross zero and ranks overlap. Joint existence, individual
magnitude and ordering precision are different questions.

\textbf{Exercise 4 --- BH calculation}

With m = 25 and q = 0.05, the first two BH thresholds are 0.002 and
0.004. Raw p-values 0.000254 and 0.0063 produce only one discovery at
these positions because the second exceeds 0.004, unless a later ordered
p-value satisfies its larger threshold.

\textbf{Exercise 5 --- Why resample common dates?}

If each portfolio resamples different dates, same-month common shocks
are destroyed. Joint score covariance and rank co-movement are then
understated.

\textbf{Exercise 6 --- Why refit \ensuremath{\mu } and \ensuremath{\tau } inside each bootstrap?}

They are estimated quantities. Holding them fixed conditions away
hyperparameter-estimation uncertainty and makes rank/score intervals too
narrow.

\textbf{Exercise 7 --- Overlap arithmetic}

Two 120-month windows 60 months apart share 60 months, so overlap is
(120\ensuremath{-}60)/120 = 50\%. They are not independent persistence observations.

\textbf{Exercise 8 --- Frozen beta}

Future alpha is evaluated relative to the risk exposure estimated from
training data. Re-estimating beta on future observations lets evaluation
data redefine the benchmark residual and introduces look-ahead.

\textbf{Exercise 9 --- Positive forward spread but no dynamic increment}

A stable cross-sectional ordering can predict future outcomes. If the
expanding-window ranking performs as well as the rolling rank, annual
re-estimation contributes no additional timing information.

\textbf{Exercise 10 --- Why is the oracle not a fair competitor?}

It uses the full sample, including evaluation periods, so it has
look-ahead. It is an upper bound used to diagnose the strength of a
fixed ordering, not an implementable strategy.

\textbf{Exercise 11 --- HAC versus GRS}

GRS has an exact F reference under IID Gaussian errors. HAC/Wald is
asymptotically robust to specified heteroskedasticity and
autocorrelation. When diagnostics reject IID Gaussian residuals, GRS
remains informative as a conventional reference but should not be
primary.

\textbf{Exercise 12 --- Source mismatch}

If local calculations reproduce exactly but local values differ from the
current official snapshot, numerical implementation is validated while
source equality is not. Report the two findings separately.

\textbf{Exercise 13 --- SFA interpretation}

A supported one-sided residual component indicates asymmetry under the
fitted distribution. It does not prove inefficiency, especially when
omitted risks, nonlinearities and tail events can also create asymmetry.

\textbf{Exercise 14 --- Design a negative control}

Simulate returns from the fitted factors with true alpha zero and
dependent residuals. The pipeline should recover near-zero posterior
centres, controlled false discoveries and no systematic forward
increment. This tests calibration rather than reproducing the empirical
result.

\textbf{Exercise 15 --- Complete the SFA square}

Starting from (\ensuremath{\varepsilon }+u)\textsuperscript{2}/\ensuremath{\sigma }\textsubscript{v}\textsuperscript{2}+u\textsuperscript{2}/\ensuremath{\sigma }\textsubscript{u}\textsuperscript{2}, collect
coefficients of u\textsuperscript{2} and u. Set
\ensuremath{\sigma }\textsubscript{*}\textsuperscript{2}=\ensuremath{\sigma }\textsubscript{u}\textsuperscript{2}\ensuremath{\sigma }\textsubscript{v}\textsuperscript{2}/(\ensuremath{\sigma }\textsubscript{u}\textsuperscript{2}+\ensuremath{\sigma }\textsubscript{v}\textsuperscript{2})
and
\ensuremath{\mu }\textsubscript{*}=\ensuremath{-}\ensuremath{\varepsilon }\ensuremath{\sigma }\textsubscript{u}\textsuperscript{2}/(\ensuremath{\sigma }\textsubscript{u}\textsuperscript{2}+\ensuremath{\sigma }\textsubscript{v}\textsuperscript{2}).
Expanding (u\ensuremath{-}\ensuremath{\mu }\textsubscript{*})\textsuperscript{2}/\ensuremath{\sigma }\textsubscript{*}\textsuperscript{2}+\ensuremath{\varepsilon }\textsuperscript{2}/\ensuremath{\sigma }\textsuperscript{2} reproduces
every coefficient, proving the identity used in both marginalisation and
conditioning.

\textbf{Exercise 16 --- Derive the SFA marginal density}

Insert the completed square into f(\ensuremath{\varepsilon },u), factor out exp{[}\ensuremath{-}\ensuremath{\varepsilon }\textsuperscript{2}/(2\ensuremath{\sigma }\textsuperscript{2}){]},
and integrate the remaining N(\ensuremath{\mu }\textsubscript{*},\ensuremath{\sigma }\textsubscript{*}\textsuperscript{2})
kernel from zero to infinity. Use
\ensuremath{\mu }\textsubscript{*}/\ensuremath{\sigma }\textsubscript{*}=\ensuremath{-}\ensuremath{\lambda }\ensuremath{\varepsilon }/\ensuremath{\sigma } and
\ensuremath{\sigma }\textsubscript{*}/(\ensuremath{\sigma }\textsubscript{u}\ensuremath{\sigma }\textsubscript{v})=1/\ensuremath{\sigma } to obtain
f(\ensuremath{\varepsilon })=(2/\ensuremath{\sigma })\ensuremath{\phi }(\ensuremath{\varepsilon }/\ensuremath{\sigma })\ensuremath{\Phi }(\ensuremath{-}\ensuremath{\lambda }\ensuremath{\varepsilon }/\ensuremath{\sigma }).

\textbf{Exercise 17 --- Verify the SFA conditional mean}

Standardise w=(u\ensuremath{-}\ensuremath{\mu }\textsubscript{*})/\ensuremath{\sigma }\textsubscript{*}. The lower bound
becomes \ensuremath{-}z\textsubscript{*}. Integration by parts gives
\ensuremath{\int }\textsubscript{\ensuremath{-}z*}\textsuperscript{\ensuremath{\infty }}w\ensuremath{\phi }(w)dw=\ensuremath{\phi }(z\textsubscript{*}).
Divide by \ensuremath{\Phi }(z\textsubscript{*}) and transform back:
E(u\textbar \ensuremath{\varepsilon })=\ensuremath{\mu }\textsubscript{*}+\ensuremath{\sigma }\textsubscript{*}\ensuremath{\phi }(z\textsubscript{*})/\ensuremath{\Phi }(z\textsubscript{*}).

\textbf{Exercise 18 --- Prove unit invariance}

Replace \ensuremath{\varepsilon },\ensuremath{\sigma }\textsubscript{u},\ensuremath{\sigma }\textsubscript{v} by a times themselves
for a\textgreater0. Then \ensuremath{\sigma },\ensuremath{\mu }\textsubscript{*},\ensuremath{\sigma }\textsubscript{*} also
multiply by a, while \ensuremath{\mu }\textsubscript{*}/\ensuremath{\sigma }, \ensuremath{\sigma }\textsubscript{*}/\ensuremath{\sigma } and
\ensuremath{\mu }\textsubscript{*}/\ensuremath{\sigma }\textsubscript{*} do not change. Hence
E{[}exp(\ensuremath{-}u/\ensuremath{\sigma })\textbar \ensuremath{\varepsilon }{]} is invariant. E{[}exp(\ensuremath{-}u)\textbar \ensuremath{\varepsilon }{]} is
not.

\textbf{Exercise 19 --- Boundary p-value}

If LR=6.635, the ordinary \ensuremath{\chi }\textsuperscript{2}\textsubscript{1} upper tail is
approximately 0.010. The 50:50 boundary-mixture p-value is approximately
0.005. If LR=0, the point mass implies p=1. Apply BH to the resulting
family of 25 SFA p-values.

\textbf{Exercise 20 --- Posterior shrinkage}

With \ensuremath{\mu }=\ensuremath{-}26, \ensuremath{\tau }=94, raw alpha a=250 and standard error s=150 bps/year, the
diagonal-model weight is 94\textsuperscript{2}/(94\textsuperscript{2}+150\textsuperscript{2})=0.282. The posterior mean is
0.282(250)+0.718(\ensuremath{-}26)=51.8 bps/year. The full-covariance answer need not
equal this because other portfolios enter through off-diagonal sampling
covariance.

